\documentclass[11pt, letterpaper]{article}
\usepackage[margin=1in]{geometry}
\usepackage{amsmath, amssymb}
\usepackage{graphicx}
\usepackage{hyperref}
\usepackage{booktabs}
\usepackage{setspace}

\title{\textbf{Deep Learning Midterm Project}\\Multi-Object Tracking on the MOT16 Dataset: A Tracking-by-Detection Approach}
\author{\textbf{Group 12} \\ Sal Nigro, Brady Maes}
\date{\today}

\begin{document}

\maketitle

\begin{abstract}
Video object tracking is a fundamental problem in computer vision, forming the basis for applications ranging from autonomous driving to intelligent video surveillance. In this report, we present a baseline multi-object tracking (MOT) system evaluated on the challenging MOT16 dataset. Our system follows the tracking-by-detection paradigm, integrating a fine-tuned Faster R-CNN for pedestrian detection, a custom Siamese Network for appearance feature extraction (Re-Identification), and the Hungarian algorithm for data association. We evaluate our pipeline on the MOT16-02 sequence using standard metrics, obtaining a Multiple Object Tracking Accuracy (MOTA) of $-15.5\%$ and an ID F1 Score (IDF1) of $3.3\%$. We discuss the architectural choices, the specific challenges encountered during implementation—such as thresholding dilemmas and embedding space collapse—and provide a detailed self-evaluation of the system's performance and computational speed. Finally, we propose concrete future directions to address the high rate of false negatives and identity switches observed in this early-stage baseline.
\end{abstract}

\section{Introduction}
Video object tracking involves the continuous localization of objects across sequential frames of a video. Unlike single-object tracking, Multi-Object Tracking (MOT) requires the algorithm to not only localize an unknown number of objects but also to maintain their distinct identities over time, even as they interact, occlude one another, or temporarily leave the camera's field of view.

\subsection{Importance and Motivation}
The motivation for studying video object tracking is rooted in its immense practical utility and its role as a precursor to high-level video understanding. 
\begin{itemize}
    \item \textbf{Autonomous Navigation:} Self-driving vehicles require accurate, real-time tracking of pedestrians, cyclists, and other vehicles to plan safe trajectories and avoid collisions.
    \item \textbf{Crowd Analytics and Surveillance:} Intelligent monitoring systems rely on MOT to estimate crowd densities, track suspicious activities, and ensure public safety in highly congested environments.
    \item \textbf{Sports Analytics:} Tracking players across a field provides statistical insights into game dynamics, player fatigue, and tactical formations.
\end{itemize}

Despite decades of research, MOT remains exceptionally challenging. Real-world video streams, such as those found in the MOT16 benchmark dataset \cite{milan2016mot16}, feature severe occlusions, erratic camera motion, varying illumination conditions, and individuals with highly similar appearances. 

\subsection{The Tracking-by-Detection Paradigm}
To tackle these challenges, the prevailing methodology in modern MOT is "Tracking-by-Detection". In this framework, the tracking problem is decoupled into two distinct stages:
1. \textbf{Detection:} An object detector locates all targets of interest in a given frame independently of past or future frames.
2. \textbf{Data Association:} A matching algorithm links the detections from the current frame to the active trajectories (tracks) established in previous frames.

Our project strictly adheres to this paradigm. By separating detection from association, we can leverage state-of-the-art Convolutional Neural Networks (CNNs) for robust localization, while utilizing specialized Re-Identification (Re-ID) networks to handle identity preservation. This report details the theoretical foundations, implementation specifics, and empirical evaluation of our custom MOT pipeline.

\section{Methodology}

Our methodology consists of three primary components: the Object Detection module, the Appearance Modeling (Re-ID) module, and the Data Association module.

\subsection{Model Introduction: Faster R-CNN}
For the detection stage, we employ Faster R-CNN (Region-based Convolutional Neural Network) \cite{ren2015faster}. Faster R-CNN is a two-stage object detector renowned for its high accuracy. 
The architecture consists of a backbone network (ResNet-50 combined with a Feature Pyramid Network, FPN) that extracts deep feature maps from the input image. These features are fed into a Region Proposal Network (RPN) which slides a small network over the feature map to output bounding box proposals and objectness scores. The proposed regions are then pooled to a fixed size using RoIAlign and passed through fully connected layers to output the final bounding box regression offsets and class probabilities. 

In our implementation, we fine-tuned a pre-trained Faster R-CNN model specifically on the MOT16 training sequences. The model was configured to predict two classes: Background (0) and Pedestrian (1), effectively filtering out irrelevant objects and focusing the network's capacity on the target class.

\subsection{Model Introduction: Siamese Network for Re-Identification}
Once pedestrians are detected, the system must determine if a newly detected person is the same individual as a previously tracked person. To achieve this, we designed a custom Siamese Network \cite{bromley1993signature} to extract appearance features.
A Siamese Network consists of two identical subnetworks that share the exact same weights and parameters. The network takes a pair of images as input, processes them through the identical convolutional layers to produce two feature embeddings, and then computes the distance or similarity between these embeddings.

Our Siamese Network accepts $24 \times 24$ grayscale crops of pedestrians. The architecture consists of multiple 2D convolutional layers followed by ReLU activations and Max Pooling. The flattened outputs of both sub-networks are concatenated and passed through a fully connected layer with a Sigmoid activation. This structure directly outputs a similarity score $S \in [0, 1]$, where $1$ indicates the images are of the same person, and $0$ indicates they are different people. The network was trained using Binary Cross-Entropy (BCE) loss on positive and negative image pairs generated dynamically from the MOT16 ground truth data.

\subsection{System Design and Data Association}
The core tracking logic bridges the detection and Re-ID models. For every frame $t$ in the video sequence:
1. The frame is passed through the Faster R-CNN model. Detections with a confidence score above a predefined threshold ($\tau_{det} = 0.7$) are retained.
2. For each retained detection, a bounding box crop is extracted from the original frame, resized to $24 \times 24$, converted to grayscale, and passed through the Siamese Network to generate an appearance feature.
3. A cost matrix $C$ of size $N \times M$ is constructed, where $N$ is the number of new detections in frame $t$, and $M$ is the number of currently active tracks. 

The elements of the cost matrix $C_{i,j}$ are computed based on the similarity score $S$ returned by the Siamese Network between the $i$-th detection and the $j$-th active track:
\begin{equation}
C_{i,j} = 1.0 - S(D_i, T_j)
\end{equation}

To assign detections to tracks optimally, we formulate this as a bipartite matching problem and solve it using the Hungarian Algorithm (Linear Sum Assignment) \cite{kuhn1955hungarian}. The algorithm minimizes the global assignment cost. We enforce a matching threshold $\tau_{match} = 0.5$; any assignment with a cost greater than this threshold is rejected. Unassigned detections initialize new tracks, while active tracks matched to new detections have their appearance features updated to adapt to visual changes over time.

\subsection{Implementation Details}
The models were implemented using PyTorch and Torchvision. 
\begin{itemize}
    \item \textbf{Dataset Parsing:} Custom PyTorch `Dataset` classes were written to parse the MOT16 `gt.txt` files, extracting frame numbers, bounding box coordinates, and unique object IDs.
    \item \textbf{Data Augmentation:} For the Siamese Network, we utilized grayscale conversion to enforce learning of structural features rather than relying purely on color, which can vary wildly under different lighting conditions.
    \item \textbf{Hardware:} Computations were accelerated using CUDA-enabled GPUs, leveraging batch processing during the fine-tuning phases.
\end{itemize}

\section{Challenges Encountered}

Implementing a multi-object tracker from scratch presents several intricate challenges. The negative MOTA score observed in our baseline evaluation highlights the inherent difficulty of the task.

\subsection{Detection Thresholding and False Negatives}
A major challenge in tracking-by-detection is the extreme sensitivity of the pipeline to the performance of the detector. If the detector misses an object (a False Negative), the tracker has no information to associate, and the track is inevitably broken. When the object is detected again in a subsequent frame, it is assigned a new ID, resulting in an ID Switch. 

To mitigate False Positives (background noise being classified as pedestrians), we applied a strict detection confidence threshold ($\tau_{det} = 0.7$). However, this introduced a vast number of False Negatives, severely impacting the continuity of our tracks. Solving this trade-off requires more exhaustive hyperparameter tuning or utilizing Non-Maximum Suppression (NMS) coupled with a lower confidence threshold.

\subsection{Siamese Network Embedding Robustness}
The Re-ID process proved highly complex. Pedestrians in the MOT16 dataset frequently wear similar clothing (e.g., dark winter coats) and are often captured in low resolution. Our Siamese Network, trained on $24 \times 24$ grayscale crops, occasionally suffered from "embedding space collapse", where it struggled to assertively differentiate between two different individuals. 

To address this during inference, we implemented strict thresholding in the Hungarian matching step ($\tau_{match} = 0.5$). If the network returned a similarity score lower than $0.5$, the matching was explicitly forbidden, forcing the creation of a new track rather than incorrectly merging two distinct individuals. While this prevented some catastrophic identity merges, it contributed heavily to the inflation of ID switches.

\subsection{Lack of Motion Modeling}
Our current data association relies entirely on appearance features. We encountered significant issues when pedestrians physically occluded one another. Because appearance alone cannot resolve which person is which when they overlap, the Hungarian algorithm often swapped their identities. 

A standard solution to this, which we identify as critical future work, is the integration of a Kalman Filter \cite{kalman1960new}. A Kalman filter predicts the future bounding box coordinates of an object based on its past velocity and trajectory. Incorporating spatial distance (e.g., Intersection over Union) alongside Re-ID similarity into the cost matrix would drastically reduce ID switches during occlusions.

\section{Self-Evaluation}

We evaluated our full tracking pipeline on the \texttt{MOT16-02} training sequence using the standard \texttt{motmetrics} Python library. The \texttt{MOT16-02} sequence consists of 600 frames recorded in a busy plaza with a moving camera and varying lighting conditions.

\subsection{Performance Analysis}
The quantitative results of our pipeline are summarized in Table \ref{tab:results}.

\begin{table}[h]
\centering
\caption{Tracking Evaluation Results on Sequence MOT16-02}
\label{tab:results}
\begin{tabular}{@{}lcccccccc@{}}
\toprule
\textbf{Sequence} & \textbf{Frames} & \textbf{MOTA} & \textbf{MOTP} & \textbf{IDF1} & \textbf{MT} & \textbf{FP} & \textbf{FN} & \textbf{ID SW} \\ \midrule
MOT16-02 & 600 & -15.5\% & 0.215 & 3.3\% & 9 & 2974 & 10670 & 6957 \\ \bottomrule
\end{tabular}
\end{tabular}
\end{table}

\textbf{Multiple Object Tracking Accuracy (MOTA):} The MOTA score aggregates three sources of error: False Positives (FP), False Negatives (FN), and ID Switches (ID SW). Our baseline achieved a MOTA of $-15.5\%$. In the MOT context, a negative MOTA indicates that the number of errors made by the tracker exceeds the total number of ground truth objects. This is primarily driven by the massive volume of False Negatives (10,670) and ID Switches (6,957). 

\textbf{ID F1 Score (IDF1):} The IDF1 score measures the ratio of correctly identified detections over the average number of ground-truth and computed detections. Our score of $3.3\%$ reflects the fragmentation of our trajectories. Due to the strict detection threshold and lack of motion modeling, tracks rarely survived for the entire duration of a pedestrian's presence in the scene.

While these metrics seem low, they are entirely expected for a foundational, appearance-only baseline implemented without advanced heuristics like Kalman filtering, track patience (keeping lost tracks alive for $N$ frames), or high-resolution feature embedding.

\subsection{Speed and Efficiency Analysis}
From a computational standpoint, our unoptimized pipeline presents a bottleneck during the data association phase. For $N$ detections and $M$ active tracks in a single frame, the system performs $N \times M$ independent inferences through the Siamese Network to populate the cost matrix. 

This quadratic complexity results in significant frame-rate drops in crowded scenes where $N$ and $M$ are large. 
\textbf{Solution applied:} To mitigate this, we constrained the maximum number of active track comparisons and aggressively purged old tracks. 
\textbf{Future Solution:} A superior architectural approach is the Joint Detection and Embedding (JDE) paradigm \cite{wang2020towards}, where a single neural network outputs both bounding boxes and Re-ID embeddings simultaneously, reducing the cost matrix generation to a simple matrix multiplication of pre-computed vectors.

\section{Conclusion}
In this project, we successfully designed and implemented an end-to-end Multi-Object Tracking pipeline using a Tracking-by-Detection methodology. By integrating a fine-tuned Faster R-CNN, a custom Siamese Re-ID network, and Hungarian matching, we established a functioning baseline on the MOT16 dataset. Through rigorous self-evaluation, we identified critical failure points in purely appearance-based data association and detection thresholding. Addressing these bottlenecks through the incorporation of motion modeling (Kalman Filters) and joint embedding architectures provides a clear roadmap for advancing this system toward state-of-the-art performance.

\begin{thebibliography}{9}

\bibitem{milan2016mot16}
A. Milan, L. Leal-Taix{\'e}, I. Reid, S. Roth, and K. Schindler, 
\textit{MOT16: A Benchmark for Multi-Object Tracking}, 
arXiv preprint arXiv:1603.00831, 2016.

\bibitem{ren2015faster}
S. Ren, K. He, R. Girshick, and J. Sun, 
\textit{Faster R-CNN: Towards Real-Time Object Detection with Region Proposal Networks}, 
Advances in Neural Information Processing Systems (NeurIPS), 2015.

\bibitem{bromley1993signature}
J. Bromley, I. Guyon, Y. LeCun, E. Säckinger, and R. Shah, 
\textit{Signature Verification using a "Siamese" Time Delay Neural Network}, 
Advances in Neural Information Processing Systems (NeurIPS), 1993.

\bibitem{kuhn1955hungarian}
H. W. Kuhn, 
\textit{The Hungarian Method for the Assignment Problem}, 
Naval Research Logistics Quarterly, 2(1-2), 83-97, 1955.

\bibitem{kalman1960new}
R. E. Kalman, 
\textit{A New Approach to Linear Filtering and Prediction Problems}, 
Journal of Basic Engineering, 1960.

\bibitem{wang2020towards}
Z. Wang, L. Zheng, Y. Liu, Y. Li, and S. Wang, 
\textit{Towards Real-Time Multi-Object Tracking}, 
European Conference on Computer Vision (ECCV), 2020.

\end{thebibliography}

\end{document}
